\documentclass{article}
\usepackage[T1]{fontenc}
\usepackage{lmodern}
\usepackage{graphicx}
\usepackage{amsmath,amssymb}
\usepackage[a4paper,margin=2cm]{geometry}
\usepackage[absolute,overlay]{textpos}
\setlength{\TPHorizModule}{1mm}
\setlength{\TPVertModule}{1mm}
\usepackage[indonesian]{babel}
\usepackage{hyperref}
\setlength{\emergencystretch}{2em}
% unit: Halaman 53

\begin{document}

% === Halaman 53 (python/native) ===

Beberapa Varian Metode LSB

1.  Sequential 

• Bit-bit pesan disembunyikan secara sekuensial pada pixel-pixel citra. 

• Misalkan ukuran pesan = 15 bit, maka urutan pixel-pixel yang digunakan 

untuk penyembunyian bit adalah:

53

2

6

3

4

5

1

7

8

10

14

11

12

13

9

15

=

\end{document}
